\section{Introduction}

% Measuring individual differences within and across species is fundamental to many areas of biology~\parencite{nelson1970outline,wagner1989biological}.
% While such comparisons have yielded significant insights in neuroscience—particularly in surveys of gross-level neuroanatomy across species~\parencite{butler2005comparative} and in studies of small, well-characterized central pattern generator circuits~\parencite{Marder2006}—it remains 
A major challenge in neuroscience is to quantitatively compare the complex, high-dimensional activity of large collections of neural systems. This challenge sits at the center of an ongoing debate about how neural codes are organized, a question that recurs at every scale of description. At the level of single neurons, responses often appear disorganized: neurons typically exhibit mixed selectivity, responding to seemingly arbitrary combinations of variables~\parencite{rigotti2013,tye2024,ostojic2024,posani2026}. At the level of brain regions, a classical view holds that function is localized, with distinct areas playing specialized roles arranged along processing hierarchies \parencite{petersen2024principles}. Recent high-throughput recordings have complicated this picture, revealing that behavioral variables can be decoded from almost anywhere in the brain~\parencite{steinmetz2019distributed,stringer2019spontaneous,musall2019, findling2025brain,angelaki2025brain,rosen2026distributed,schartner2025, bondy2025}, a result interpreted as evidence that population codes are broadly distributed and largely undifferentiated across regions \parencite{hayden2026rethinking,fine2021whole,Lloreda2025}.
\textcite{posani2026} argue against this reading by showing that a neuron's area can be decoded from its selectivity profile.
But arguments on both sides characterize each area on its own---through single-variable decoding or summaries of single-neuron tuning---rather than comparing areas to one another.
Furthermore, most analyses pool neurons across subjects, assuming that variability from one animal to the next is idiosyncratic noise rather than structured signal.

Part of these debates could be resolved by moving from single-neuron tuning and single-variable decodability to a fully geometrical picture, i.e. how the full set of behavioral variables is jointly arranged and how these arrangements differ within and across scales. Indeed, many have argued that this arrangement, often called neural geometry, directly reflects a system's computation (\cite{bernardi,courellis2024,nogueira2026, fascianelli2024neural, pagan2025, HarveyWilliamsDecoding}), and numerous measures have been proposed to quantitatively compare these geometries across systems, beginning with Representational Similarity Analysis (RSA; \cite{Kriegeskorte2008_frontiers}) and its many variants \parencite[see][for a review]{kriegeskorte2019peeling}, alongside Canonical Correlation Analysis \parencite[CCA;][]{Raghu2017,gallego2020long} and hyperalignment \parencite{haxby2011common,haxby2020hyperalignment}. %, as well as regression- and classification-based comparisons of representational geometry \parencite{bernardi,conwell2022can,fascianelli2024neural,muzellec2026reverse}.

Crucially, scaling up these comparative analyses to systematically analyze large collections of subjects, sessions, and brain regions (\textbf{Fig.} \ref{fig:schematic}a) remains a challenge.
For instance, one might ask whether the geometries fall into clusters---i.e., small groups of subjects or brain regions whose population activity sits closer to one each other than to the rest. Similar questions have been asked at the level of single-neuron tuning proprieties (\cite{hirokawa2019,posani2026}), but only rarely for entire recordings, and even then for small numbers of subjects \parencite{safaie2023preserved,Chen2021,gallego2020long}. 
One might also ask whether a particular neural geometry is behaviorally relevant---e.g., by demonstrating that a subject's behavioral profile can be inferred from other subjects with similar geometries. 
However, these analyses presuppose that the systems occupy a common space, each recording holding a position within it and distances across systems reflecting their similarity.

Here, we leverage Procrustes analysis \parencite{Gower2004,Williams2021,pospisil2023estimating,harvey2024duality,khosla2024soft} as a principled approach to place neural geometries within a common \textit{shape space} over which clustering, prediction, and low-dimensional visualization can be performed rigorously.
% The mathematical foundations of this approach are classical \parencite{kendall2009shape}, and their extension to neural geometries  has been extensively detailed in recent methodological papers \parencite{Williams2021,pospisil2023estimating,harvey2024duality,khosla2024soft}.
% We apply this framework to show that neural population geometry is systematically structured across both regions and individuals in ways that single-neuron analyses and decoding performance scores fail to capture. 
Applying this framework, we find that neural population geometry is systematically structured at every scale we examined. In simulations, we first show that structure at the neuron level and structure across systems can be varied independently, so that clustering neurons pooled across regions is blind to how the regions themselves are organised (\textbf{Fig.}~1). Applying shape metrics to multi-regional recordings in mice \parencite{angelaki2025brain} and monkeys \parencite{Siegel2015-cz}, we find that although task variables are decodable nearly everywhere, regional geometry recovers the anatomical hierarchy without supervision \parencite{posani2026}. Importantly, areas form a continuum rather than discrete clusters (\textbf{Fig.}~2). Across individuals, head-direction populations share a common ring manifold that supports zero-shot decoding across animals, but each animal carries reliable, idiosyncratic distortions organised by tuning statistics (\textbf{Fig.}~3). Finally, the same pipeline that finds no clusters across cortical areas separates wild-type from mutant mice carrying autism-associated mutations, indistinguishable at the single-neuron level \parencite{noel2025}, and links each animal's geometry to its behavioural strategy (\textbf{Fig.}~4). 

% \hl{update results summary}
% Applying this framework across dozens of brain regions in mice and non-human primates, we recover a hierarchical organization of brain regions in a fully unsupervised manner, even though individual task variables are decodable from nearly everywhere in these brains. In the mouse head direction system, we uncover systematic animal-to-animal differences in head direction coding manifolds, trace these differences to tuning width via regression in shape space, and show that all subjects nonetheless share a common ring-shaped manifold, enabling decoding of head direction in a held-out animal from other animals' data alone \parencite{safaie2023preserved}. Further, in a large cohort of mice performing a decision-making task, we show that individual differences in behavior are mirrored by individual differences in neural population geometry in most considered regions. Finally, we find that genetically defined differences between animals, invisible at the single-neuron level~\parencite{noel2025}, are reflected in population geometry.

% Together, these results show that neural population geometry is structured across both regions and individuals, organized by anatomy across brain areas and by behavior and genotype across individuals.